\section{A GPU interior-point solver}\label{sec:algorithm}

\subsection{Overview}

SPRINT is a primal--dual interior-point solver for problems of the
form~\eqref{eq:sdp}. It uses a homogeneous self-dual
embedding~\cite{ye1994homogeneous}, Mehrotra predictor--corrector
steps~\cite{mehrotra1992implementation} and double-precision arithmetic. The
outer method owns the scaling, regularization, step lengths and termination
tests. A pluggable linear backend owns symbolic analysis, numerical
factorization and solves for the current Schur operator, so the same IPM policy
can run with CPU sparse direct solvers, full-system cuDSS, or the structured
backend described here. \Cref{fig:pipeline} summarizes the structured backend.
Its main design principle follows from \cref{sec:hierarchy,sec:complexity}:
keep local numerical work local, and keep only the temporal coupling that must
remain global. \todo{state the SPRINT code availability / repository link.}

\begin{figure}[htbp]
\centering
\begin{tikzpicture}[
  node distance=5mm and 5mm,
  box/.style={draw=paperblue,rounded corners=2pt,align=center,font=\footnotesize,
              minimum height=11mm,text width=22mm,fill=paperblue!6},
  setup/.style={box,draw=paperteal,fill=paperteal!8},
  fb/.style={box,draw=paperorange,fill=paperorange!10},
  arr/.style={-{Stealth[length=2mm]},thick,paperblue}]
  \node[box] (c) {Selector Schur\\assembly};
  \node[box,right=of c] (d) {Batched local\\partial Cholesky};
  \node[box,right=of d] (e) {Reduced solve:\\dense / PCR / cuDSS};
  \node[box,right=of e] (f) {Condense,\\recover, rank-one\\correction};
  \node[box,right=of f] (g) {Full-operator\\residual gate};
  \node[setup,above=11mm of c] (a) {SOS SDP with\\stage metadata};
  \node[setup,right=of a] (b) {Certified\\symbolic plan};
  \node[box,above=11mm of g,fill=white] (ipm) {Accepted\\Newton direction};
  \node[fb,below=9mm of g] (h) {Lazy fallback:\\full cuDSS};
  \node[draw=papergray,densely dashed,rounded corners,inner sep=2.5mm,
        fit=(c)(d)(e)(f)(g)] (loop) {};
  \node[font=\footnotesize,text=papergray,anchor=north west] at (loop.south west)
        {repeated for each operator generation (IPM iteration)};
  \draw[arr] (a) -- (b);
  \draw[arr] (b.south) -- (b.south |- loop.north);
  \draw[arr] (c) -- (d);
  \draw[arr] (d) -- (e);
  \draw[arr] (e) -- (f);
  \draw[arr] (f) -- (g);
  \draw[arr] (g) -- node[right,font=\footnotesize]{accept} (ipm);
  \draw[arr,paperorange] (g) -- node[right,font=\footnotesize]{reject} (h);
  \node[setup,fill=none,draw=none,text width=34mm,right=2mm of b,
        font=\footnotesize\itshape,text=papergray]
        {once per problem: separators, pruning, coarsening, certificate, supernodal plans};
\end{tikzpicture}
\caption{Structure of the SPRINT stage backend. Setup (teal) runs once per
problem. Each IPM operator generation assembles the Schur matrix, factors the
interior components in batches, solves a small reduced system, and accepts a
direction only after a residual check against the full operator; on rejection,
a lazily created full-system cuDSS factorization of the same operator is used.}
\label{fig:pipeline}
\end{figure}

\subsection{Symbolic preparation}

The customized SPOT exporter attaches stage metadata to PSD blocks, free
variables and coefficient rows, while leaving the SDP coefficients unchanged.
The metadata is used only to \emph{propose} interfaces; the assembled
conservative Schur pattern is the final authority. Initial separators come from
stage-local connectivity and vertex cuts. Carry pruning and coarsening then
propose smaller retained sets, and every proposal is accepted only after the
certificate of \cref{def:certificate} has been checked on the actual pattern.
When the last physical cut disappears, its tail joins an interior component and
the remaining groups are reindexed consecutively. CHOLMOD~\cite{chen2008cholmod}
is used for symbolic ordering only; the numerical factorization runs on arrays
owned by SPRINT. \Cref{alg:symbolic} summarizes the procedure.

\begin{algorithm}[htbp]
\caption{Symbolic preparation (once per problem)}\label{alg:symbolic}
\small
\begin{algorithmic}[1]
\Require SDP coefficient maps, conservative Schur pattern $G$, stage metadata, border $B$
\Ensure a certified structured symbolic plan, or a general-solver route
\State validate row indices, stage metadata and the complete Schur pattern
\State propose stage separators; optionally apply carry pruning and coarsening
\State form disjoint groups $R_1,\ldots,R_K$ with compact indices
\State find the components $I_j$ of $G[V\setminus R]$
\If{a retained edge or a component boundary violates \cref{def:certificate}}
  \State \Return general-solver route
\EndIf
\State build local interior/front patterns and scatter maps from the global Schur storage
\State order each interior with its retained front last; build supernodal plans
\State group components with compatible patterns; build packing and execution schedules
\State build the reduced chain/border pattern and select its solver route
\State allocate reusable buffers and publish the plan
\end{algorithmic}
\end{algorithm}

The linear-time certificate check does not bound every step of
\cref{alg:symbolic}: separator search, ordering, pattern grouping and the
current indexing structures have their own costs (\cref{app:implementation}).

\subsection{Assembly and local partial factorization}\label{sec:local-factor}

At each iterate, the solver computes the cone scalings and assembles the
numerical Schur matrix, choosing per block between sparse-coordinate selection
(\cref{prop:selector}) and a dense block path. The current implementation
assembles the global sparse matrix before packing local panels.

For component $j$, only the $A_j$ variables of the local augmented matrix
$\bigl[\begin{smallmatrix}A_j&G_j\\G_j^\top&0\end{smallmatrix}\bigr]$
are eliminated, so its retained block becomes $-G_j^\top A_j^{-1}G_j$. The
augmented matrix need not be positive definite; positive definiteness is
required only of the interior pivots and of the assembled reduced matrix. With an
optional positive diagonal equilibration $\Lambda_j$,
\begin{equation}\label{eq:local-factor}
\Lambda_jA_j\Lambda_j=L_jL_j^\top,\qquad C_j=L_j^{-1}\Lambda_jG_j,
\end{equation}
the local update is $-C_j^\top C_j$. For a right-hand side, with
$q_j=L_j^{-1}\Lambda_jf_j$, condensation and recovery are
\begin{equation}\label{eq:local-solve}
\widehat f_R=f_R-\sum_jC_j^\top q_j,\qquad x_j=\Lambda_jL_j^{-\top}(q_j-C_jx_R).
\end{equation}
These formulas describe the panel operations; no dense matrix holding the $C_j$
is formed. Equilibration is a change of coordinates and adds no regularization.

\subsection{Repeated patterns and scheduling}

Most stage interiors share the same sparsity pattern. The solver groups such
components, optionally taking a union of compatible coupling patterns with
explicit zeros, so that each group shares one symbolic plan while numerical
values and factors remain per instance. Within a plan, an update task runs once
its source panel is ready. Independent components and nonconflicting task
groups run concurrently, overlapping writes follow an explicit order, and
repeated schedules are captured as CUDA graphs to reduce launch overhead. The
retained front is an update destination throughout and is never factored as a
local pivot. Union patterns and relaxed supernodes add structural zeros, so the
cost~\eqref{eq:panel} is evaluated on the padded plans.

\subsection{Reduced system, recovery and the numerical contract}

After the local fronts are computed, the solver forms
$S=D+\sum_jF_j$ with $F_j=-G_j^\top A_j^{-1}G_j$, adding contributions to a
shared entry in a fixed order. The reduced system is solved by dense Cholesky
when it is small, by block cyclic reduction (PCR) when the certified chain and
the backend limits allow it, and by sparse cuDSS otherwise. A dense route is
consistent with the asymptotic bounds only under a size threshold independent
of~$N$, and a sparse route inherits the complexity of its own ordering. After the
retained variables are solved, the backward pass recovers the interior variables
from the same local factors, and the rank-one correction~\eqref{eq:sm} is
applied. \Cref{alg:factor,alg:solve} state the dependencies and the acceptance
contract.

\begin{algorithm}[htbp]
\caption{Numerical factorization for one operator generation}\label{alg:factor}
\small
\begin{algorithmic}[1]
\Require symbolic plan, current cone scalings, shift $\delta$
\Ensure published factors and rank-one preparation, or a reported failure
\State assemble $H_0$ and the rank-one vector $u$; form $H_\delta=H_0+\delta I$ once
\State pack local panels from $H_\delta$ and apply consistent equilibration
\State factor the interior panels (batched by template) and accumulate the retained fronts
\State check every local factor status and finite-value flag
\State assemble and factor the reduced operator on the selected route
\State if $u\ne0$, compute $v=H_\delta^{-1}u$ and $1+u^\top v$
\State validate and publish the factor generation
\Statex \emph{On failure:} invalidate all dependent state; a fallback may factor the
same operator; if the IPM changes $\delta$, restart for the new generation.
\end{algorithmic}
\end{algorithm}

\begin{algorithm}[htbp]
\caption{Checked solve with a published factor generation}\label{alg:solve}
\small
\begin{algorithmic}[1]
\Require published factors, right-hand side $f$, tolerance, refinement budget
\Ensure an accepted solution of the full operator $\Hfull x=f$, or failure
\State condense $f$ through the local forward factors
\State solve the reduced system for $x_R$; recover all interior variables
\State apply the rank-one correction of the same generation
\State compute the residual against the original full operator
\If{statuses, finite values and residual pass} \State \Return $x$ \EndIf
\State refine within the budget, or reject the route and solve the same
  operator with the lazily created fallback; report failure if no route succeeds
\end{algorithmic}
\end{algorithm}

\paragraph{Integration with the IPM.}
An accepted factorization serves several right-hand sides, but the corrector
right-hand side depends on the predictor and cannot be prepared in advance. Reuse
is therefore organized around one operator generation and a sequence of solves.
A rejected factor or a changed shift invalidates its rank-one preparation,
reduced factor and saved recovery data, and an affine direction computed for one
operator is never combined with a corrector from an incompatible factor state.

\subsection{Numerical acceptance}\label{sec:acceptance}

\Cref{thm:chain,thm:main} hold in exact arithmetic. In floating point, a valid
certificate does not prevent loss of accuracy, so the residual is computed with
the original assembled operator after recovery and the rank-one correction; a
reduced residual alone cannot detect an inaccurate condensation or recovery. For
$M=\Hfull\succ0$, $f\ne0$ and a computed $\widehat x$, let
$r_f=f-M\widehat x$ and $\eta_2=\|r_f\|_2/\|f\|_2$.

\begin{proposition}[residual interpretation]\label{prop:residual}
If $x=M^{-1}f$, then $\|\widehat x-x\|_M=\|r_f\|_{M^{-1}}$ and
$\|\widehat x-x\|_2/\|x\|_2\le\kappa_2(M)\,\eta_2$.
\end{proposition}

\begin{proof}
Since $\widehat x-x=-M^{-1}r_f$, the squared energy norm is $r_f^\top M^{-1}r_f$.
Also $\|\widehat x-x\|_2\le\|M^{-1}\|_2\|r_f\|_2$ and $\|f\|_2\le\|M\|_2\|x\|_2$.
\end{proof}

\begin{proposition}[effect of the shift]\label{prop:shift}
If $M_0=H_0+uu^\top\succ0$, $\delta\ge0$, $x_0=M_0^{-1}f$ and
$x_\delta=(M_0+\delta I)^{-1}f$, then
$\|x_\delta-x_0\|_2/\|x_0\|_2\le\delta/(\lambda_{\min}(M_0)+\delta)$.
\end{proposition}

\begin{proof}
The resolvent identity gives $x_\delta-x_0=-\delta(M_0+\delta I)^{-1}x_0$, and the
multiplying matrix has spectral norm $\delta/(\lambda_{\min}(M_0)+\delta)$.
\end{proof}

A small residual therefore does not imply a small forward error independently of
conditioning, and it does not certify feasibility or optimality for the SDP;
those remain the job of the outer termination test and of the independent
checks in \cref{sec:experiments}. Some stabilized routes factor an auxiliary
matrix $H_\delta+D_{\mathrm{priv}}$ and use it as a preconditioner for the target
system. Acceptance then uses the target operator with $D_{\mathrm{priv}}$
excluded, and the solver reports failure when the contract cannot be met.

\subsection{Choosing a hierarchy}\label{sec:select}

For a candidate retention and coarsening configuration $\theta$,
\cref{thm:main} suggests the structural estimate
\begin{equation}\label{eq:cost-model}
\widehat W(\theta)=\Phi_2(\theta)+\mathcal B_2(\theta)+K(\theta)p(\theta)^3
+K(\theta)p(\theta)^2g+K(\theta)p(\theta)g^2+g^3,
\end{equation}
in which~\eqref{eq:ragged} and~\eqref{eq:panel} can replace the conservative
terms. Coarsening usually reduces the number of groups but enlarges the
components, and pruning can shrink interfaces but increase local fill, so
separator cardinality alone is an incomplete objective. The current
implementation uses structural gates and empirically tuned thresholds. Selecting
the hierarchy automatically by minimizing~\eqref{eq:cost-model}, extended with
launch and bandwidth terms, is future work.
